\documentclass[a4paper]{article}
\usepackage{amsmath,amssymb,amsthm}
\usepackage{fullpage}
\usepackage{hyperref}

\newtheorem{theorem}{Theorem}[section]
\newtheorem{proposition}[theorem]{Proposition}
\newtheorem{lemma}[theorem]{Lemma}
\newtheorem{remark}[theorem]{Remark}
\DeclareMathOperator{\tr}{tr}
\DeclareMathOperator{\ord}{ord}
\DeclareMathOperator{\Nm}{N}
\newcommand{\F}{\mathbb F}
\newcommand{\Q}{\mathbb Q}
\newcommand{\Z}{\mathbb Z}
\setlength{\parindent}{0pt}
\setlength{\parskip}{0.7em}

\title{Conductor Descent for Prime-Power Lights Out:\\
Further Attempts and Their Exact Limitations}
\author{}
\date{September 5, 2026}

\begin{document}
\maketitle

\begin{abstract}
This addendum investigates the general identity
$d(p^k-1)=d(p-1)$ rather than additional prime-specific
zero-nullity families.  It does not prove the general identity or
give a counterexample.
We prove a precise obstruction to descent by powering in the regular-map
operations, an exclusion of one projected mixed-collision phase, and an
exact cyclotomic norm congruence.  Each result is separated from the
stronger, still unproved assertion that would be needed for conductor
descent.
\end{abstract}

\section{The target}

Let $\mu_N\subset\overline{\F}_2^\times$ denote the $N$th roots of unity.
For $p>5$, the outstanding assertion is that
\begin{equation}\label{eq:sum}
 1+u+u^{-1}+v+v^{-1}=0,\qquad u,v\in\mu_{p^j},
\end{equation}
forces both coordinate orders to divide $p$.
The known proof reduces the problem to
$j\leq c_p=v_p(2^{\ord_p(2)}-1)$.
The companion note \cite{endpoint} covers the additional order classes
$3\mid\ord_p(2)$ and $v_2(\ord_p(2))=2$.
Neither reduction disposes of all possible Wieferich primes.

The substitution
\[
 u=xy,\qquad v=x/y
\]
is invertible on $\mu_{p^j}^2$, because its exponent matrix has
determinant $-2$.  It converts \eqref{eq:sum} into
\begin{equation}\label{eq:product}
 (x+x^{-1})(y+y^{-1})=1.
\end{equation}
It preserves the maximum coordinate order, but need not preserve the
individual orders.  In particular, ``mixed'' in the additive coordinates
and ``mixed'' in the product coordinates are not interchangeable
notions.

\section{The group-theoretic route}

\begin{proposition}[Exact group formulation]\label{prop:group}
For $p>5$, a higher-level solution is equivalent to a configuration in
some $\mathrm{SL}_2(\F_{2^h})=\mathrm{PSL}_2(\F_{2^h})$ with
\[
 |B|=2,\qquad |A^2B|=3,\qquad
 |A|,\ |AB|\text{ powers of }p,
 \qquad \max(|A|,|AB|)\geq p^2.
\]
\end{proposition}

\begin{proof}
Given \eqref{eq:product}, put $a=x+x^{-1}$.  Then
$a\ne0,1$: the exceptional values would give $x$ order $1$ or $3$.
Set
\[
 A=\begin{pmatrix}a&1\\1&0\end{pmatrix},
 \qquad B=\begin{pmatrix}1&a+a^{-1}\\0&1\end{pmatrix}.
\]
Both determinants are one, $B$ is a nonidentity involution, and
\[
 \tr(A)=a,\qquad \tr(AB)=a^{-1}.
\]
The eigenvalues of $A$ are $x,x^{-1}$ and those of $AB$ are
$y,y^{-1}$, so their orders are exactly the coordinate orders.
Cayley--Hamilton gives
\[
 A^2=aA+I,\qquad
 \tr(A^2B)=a\tr(AB)+\tr(B)=1.
\]
A determinant-one matrix of trace one in characteristic two has
characteristic polynomial $T^2+T+1$ and exact order three.

Conversely, the prescribed orders give $\tr(B)=0$ and
$\tr(A^2B)=1$.  Cayley--Hamilton therefore implies
$\tr(A)\tr(AB)=1$.
The eigenvalues of the two odd-order matrices recover
\eqref{eq:product}, with the required maximum coordinate order.
\end{proof}

Thus the trace-product formulation is not irrelevant to the original
trace-sum problem.  The coordinate change must be taken into account
when interpreting the discussion in \cite{plateau}.

The generating pair describes an orientably regular map with face
length $|AB|$ and vertex valency $|A|$.
For $\gcd(j,|A|)=1$, Wilson's hole operation $H_j$ replaces
$(A,B)$ by $(A^j,B)$; its face length is $|A^jB|$ \cite{holes}.
The next proposition tests the most direct proposed descent.
It also allows noninvertible exponents, which are not the usual
Wilson operations but could otherwise have reduced the vertex order.

\begin{proposition}[No locus-preserving powering descent]\label{prop:holing}
Let $n=|A|$ be odd and let $A,B$ be as in the preceding construction.
For every positive integer $j$,
\[
 |A^{2j}B|=3\quad\Longleftrightarrow\quad j\equiv\pm1\pmod n.
\]
Consequently a powering step that retains the defining order-three
condition preserves both the vertex and face orders.
\end{proposition}

\begin{proof}
Let $D_0(X)=0$, $D_1(X)=X$ and
$D_{j+1}(X)=XD_j(X)+D_{j-1}(X)$ over $\F_2$.
The trace recurrence and its initial values give
\[
 \tr(A^jB)=\frac{D_j(a)}{a^2}.
\]
Since $D_{2j}(a)=D_j(a)^2$,
\[
 |A^{2j}B|=3
 \quad\Longleftrightarrow\quad D_j(a)^2=a^2
 \quad\Longleftrightarrow\quad D_j(a)=a.
\]
Writing the eigenvalues of $A$ as $x,x^{-1}$, the last equality is
$x^j+x^{-j}=x+x^{-1}$.
The quadratic $T^2+aT+1$ has exactly the roots $x,x^{-1}$.
It follows that $x^j=x$ or $x^{-1}$, equivalently
$j\equiv1$ or $-1\pmod n$.
In this case $\tr(A^jB)=a^{-1}$ as well, so the face order is unchanged.
\end{proof}

Replacing $(A,B)$ by $(AB,B)$ swaps the two odd orders and retains
the trace-product condition.  Group automorphisms preserve orders.
Therefore a sequence of these operations whose intermediate
configurations all retain the order-three condition cannot perform
conductor descent.  This does not rule out a sequence which leaves
that locus and returns using an additional argument.

There is also an entirely group-theoretic explanation for a tempting
but unhelpful observation.  Put $R=A^2B$.  From $B^2=R^3=1$,
\[
 A^4B=R A^{-2}R^{-1}.
\]
Indeed, with $U=A^2=RB$,
$RU^{-1}R^{-1}=RBR=UR=U^2B$.
Thus $|A^4B|=|A|$ for every odd-order $A$.
The equal-order map obtained after two $H_2$ operations is automatic;
it does not impose a restriction on prime-power exponents.

\section{A lower-algebra test for mixed collisions}

Write $F_0=0$, $F_1=1$, $F_{n+1}=XF_n+F_{n-1}$, and, for odd $n$,
\[
 H_n=F_{(n-1)/2}+F_{(n+1)/2}.
\]
Then $D_n=XH_n^2$, and $H_n$ is square-free with roots the nonzero
traces of nonidentity $n$th roots of unity.
Fix $p>3$, put $m=(p-1)/2$, and write
\[
 H=H_p,\qquad H^*(X)=X^mH(1/X),\qquad R(X)=D_p(1/X).
\]
Denominators in the following statements are invertible modulo $H$,
because $H(0)=1$.

\begin{lemma}[Exact reciprocal mixed criterion]\label{lem:mixed}
For $\beta$ a root of $H$, the reciprocal $\beta^{-1}$ is the
trace of a primitive $p^2$th root if and only if
\[
 H(R(\beta))=0.
\]
\end{lemma}

\begin{proof}
Choose $u$ satisfying $u+u^{-1}=\beta^{-1}$.
Then $u^p+u^{-p}=R(\beta)$.
If the latter is a root of $H$, the two solutions of its reciprocal
quadratic are primitive $p$th roots, so $u^p$ has exact order $p$.
Hence $u$ has exact order $p^2$.
The converse follows from the same trace identity.
In particular $R(\beta)=0$ is an old collision, not a mixed one:
$H(0)=1$ already excludes it from this criterion.
\end{proof}

\begin{proposition}[Exclusion of the fixed projected phase]
\label{prop:fixed-phase}
Define
\[
 Q_p(X)=H^*(X)+X^{m+1}+XH(X).
\]
Then
\[
 \gcd(H,R-X)=\gcd(H,Q_p),\qquad
 \deg Q_p<m,\qquad Q_p(0)=1.
\]
In particular, if $H_p$ is irreducible, then
$D_p(\beta^{-1})=\beta$ is impossible for a root $\beta$ of $H_p$.
\end{proposition}

\begin{proof}
The Dickson identity gives
\[
 R(X)-X=X^{-p}\bigl(H^*(X)+X^{m+1}\bigr)^2.
\]
Clearing the invertible denominator and removing the square does not
change the gcd with the square-free polynomial $H$.
Adding $XH$ does not change it either.

The two highest coefficients of $H$ are $1,1$, and $H(0)=1$.
Thus the degree-$m+1$ and degree-$m$ terms of $Q_p$ cancel, while its
constant term is one.
An irreducible degree-$m$ polynomial cannot share a root with this
nonzero lower-degree polynomial.
\end{proof}

This is an exclusion of one branch, not of all mixed collisions.
If $H$ is irreducible and one mixed collision exists, Frobenius
equivariance and transitivity imply
\[
 R(\beta)=\beta^{2^k}
\]
for a fixed phase $k$ on the entire root orbit.
Proposition~\ref{prop:fixed-phase} excludes $k=0$ only.
For every odd $a$,
\[
 \gcd(H,R-D_a)=
 \gcd\!\left(H,H^*+X^{m+1}H_a\right),
\]
by the same squared-numerator argument.
Choosing an odd representative $a\equiv\pm2^k\pmod p$ reduces
the other phases to explicit polynomials, but does not prove their
coprimality.  If $H$ is reducible, even the same-orbit phase
assumption is unjustified: one lower orbit may map to another.

\begin{proposition}[Only one candidate phase when $H_p$ is irreducible]
\label{prop:one-phase}
Assume $H=H_p$ is irreducible.  Define
\[
 \mathcal Q=\frac{H^*}{H}\pmod{X^{m+1}},\qquad
 t=\deg\mathcal Q,\qquad
 \mathcal B=\frac{H^*+\mathcal QH}{X^{m+1}},
\]
where the representative $\mathcal Q$ has degree at most $m$.
Then a reciprocal mixed collision of orders $p^2,p$ exists if and only if
\[
 t\geq1\quad\text{and}\quad
 \mathcal B=H_{2t-1}.
\]
\end{proposition}

\begin{proof}
The truncated inverse exists because $H(0)=1$, and the defining
congruence makes $\mathcal B$ a polynomial.
Since $H$ is irreducible, its roots form one Frobenius orbit of size
$m$.  The signed powers of two therefore exhaust the nonzero residue
classes modulo $p$ up to sign.
By Lemma~\ref{lem:mixed}, a mixed collision is equivalent to
$R(\beta)=D_a(\beta)$ for some odd $a$ with $1\leq a\leq p-2$.
The preceding squared-numerator identity makes this equivalent to
\[
 H^*+X^{m+1}H_a=CH
\]
for a polynomial $C$.
Comparing leading degrees gives $\deg C=(a+1)/2\leq m$.
Reducing modulo $X^{m+1}$ then gives $C=\mathcal Q$.
It follows that $a=2t-1$ and $\mathcal B=H_a$.

Conversely, the displayed equality with $t\geq1$ supplies such an odd
$a=2t-1\leq p-2$, so $R$ maps a root of $H$ to a root of $H$.
Lemma~\ref{lem:mixed} supplies the mixed collision.
If $t=0$, no quotient $C$ of the required positive degree exists.
\end{proof}

This reduces the possible phases to one explicit identity in the lower
polynomial algebra.  It is not a proof that the identity always fails.
For example, at $p=11$,
\[
 H=X^5+X^4+X^2+X+1,\qquad
 \mathcal Q=X^5+X^4+X^2+1,\qquad
 \mathcal B=X^4+X^2+1\ne H_9.
\]
Simple coefficient tests do not uniformly decide the identity:
at $p=1931$, one gets $t=964$, and $\mathcal B$ agrees with
$H_{1927}$ in degrees zero through eight before differing in degree
nine.  No argument excluding the equality for arbitrary $p$ was found.

\begin{remark}[A countermodel to dropping prime support]
Reciprocal mixed descent is false for general odd periods.
The polynomial $H_7=X^3+X^2+1$ has reciprocal
$X^3+X+1$, and
\[
 H_9=(X+1)(X^3+X+1).
\]
Thus the reciprocal of a primitive $7$-trace is a primitive $9$-trace.
Both orders divide $21^2$, while $9$ does not divide $21$.
This is not a prime-power counterexample; it shows that an argument
using only oddness and $M\mid M^2$ cannot suffice.
\end{remark}

\begin{remark}[A failed auxiliary cubic obstruction]
For the remaining irreducible cases with
$\ord_p(2)=p-1$, $m=(p-1)/2$ odd and $3\nmid m$, one might hope that
every reciprocal lift of a $p$-trace has order divisible by three.
That would exclude a $p$-primary lift, but the auxiliary assertion is
false.

Take $p=107$ and
\[
 K=\F_2[z]/(1+z+\cdots+z^{106}),\qquad \beta=z+z^{-1}.
\]
Encode an element by placing its coefficient of $z^i$ in bit $i$.
The element
\[
 u=\texttt{0x21500a7574a334aa348c8366ce2}
\]
then satisfies, by exact polynomial reduction,
\[
 (u+u^{-1})\beta=1,\qquad
 \ord(u)=3002399751580331
       =107\cdot28059810762433.
\]
This order is not divisible by three.
It is also not a power of $107$, so this example refutes only the
auxiliary assertion, not the Lights Out conjecture.
\end{remark}

\section{Exact local valuations do not force global descent}

\begin{lemma}[Simple two-adic intersection]\label{lem:carry}
Let $u,v$ satisfy \eqref{eq:sum}, and let $U,V$ be their Teichmuller
lifts in an unramified extension of $\Q_2$.
Put
\[
 S=1+U+U^{-1}+V+V^{-1}.
\]
If the coordinate orders are powers of $p>5$, then $v_2(S)=1$.
\end{lemma}

\begin{proof}
Let $\sigma$ be the unramified Frobenius, and put
$A=U+U^{-1}$, $B=V+V^{-1}$.
Direct expansion gives
\[
 S^2-\sigma(S)=2(S+1+AB).
\]
Write $S=2W$, as its reduction vanishes.
Dividing by two and reducing modulo two yields
\[
 \overline W^{\,2}=1+\alpha\beta
 =1+\alpha+\alpha^2,\qquad
 \alpha=u+u^{-1},\quad\beta=\alpha+1.
\]
If the right side vanished, multiplying by $u^2$ would give
$u^4+u^3+u^2+u+1=0$, so $u$ would have order $5$.
Thus $\overline W\ne0$ and $v_2(S)=1$.
\end{proof}

\begin{proposition}[Norm congruence]\label{prop:norm}
Let $N=p^j$, $p>5$, $L=\Q(\zeta_N+\zeta_N^{-1})$, and
\[
 S=1+\zeta_N^a+\zeta_N^{-a}+\zeta_N^b+\zeta_N^{-b},
 \qquad M=\Nm_{L/\Q}(S).
\]
Then
\begin{equation}\label{eq:norm-congruence}
 M\equiv\left(\frac5p\right)\pmod{p^j}.
\end{equation}
If $h=\min\{r>0:2^r\equiv\pm1\pmod{p^j}\}$ and exactly $t$ primes
of $L$ above two divide $S$, then
\[
 v_2(M)=ht.
\]
\end{proposition}

\begin{proof}
At the unique prime above $p$, every conjugate of $S$ is congruent
to $5$, so $p\nmid M$ and
\[
 M\equiv5^{[L:\Q]}\equiv(5/p)\pmod p.
\]
For every rational prime $\ell\ne p$, the residue degree $e$ of a
prime of $L$ above $\ell$ is the least positive integer with
$\ell^e\equiv\pm1\pmod N$.
Hence the norm of each prime-ideal factor of $(S)$ is $\pm1$
modulo $N$.  Their product gives $M\equiv\pm1\pmod N$, with its sign
uniquely selected by the congruence modulo $p$.
This proves \eqref{eq:norm-congruence}.
The second assertion follows from Lemma~\ref{lem:carry} at every
prime above two and the prime-ideal norm formula.
\end{proof}

On a Wieferich plateau, $2^h\equiv1$ or $-1\pmod{p^j}$.
Therefore the exact two-part $2^{ht}$ is already compatible with
the norm congruence.  The odd part and the sign of $M$ are not
controlled by these statements.

Nor do arbitrary global congruences repair this issue.
For any prescribed subset of the primes above two, the Chinese
remainder theorem produces an algebraic integer divisible simply at
exactly those primes, while agreeing with $S$ to any prescribed
finite $p$-adic precision.
Such an integer need not have the special five-term shape or its
small height.  Thus this observation is not a counterexample; it
pinpoints the information missing from an argument using only local
valuation and norm congruences.

\section{What this investigation has and has not established}

The direct locus-preserving powering descent cannot work.
The global norm calculation is exact but consistent with a higher
conductor.  Under an irreducibility hypothesis, the mixed-collision argument
eliminates the fixed projected phase and reduces all remaining phases
to one explicit polynomial identity.  It does not prove that the
identity fails, or handle collisions between distinct lower orbits.

Consequently none of these three attempts proves the all-prime
identity, and no prime-power counterexample has been obtained.
In particular, the equivalent group formulation must not be promoted
to a group-theoretic exponent bound without an additional argument.

\begin{thebibliography}{9}
\bibitem{endpoint}
\emph{Two Trace--Norm Theorems for Prime-Power Lights Out Stabilization},
\path|tex\prime_power_endpoint\prime_power_endpoint.tex|.

\bibitem{plateau}
W.~Boyles, \emph{Reciprocal Five-Term Relations among $p^j$-th Roots
of Unity in Characteristic Two: Reductions, Obstructions, and the
Wieferich Plateau},
\path|tex\five_term_plateau\five_term_plateau.tex|.

\bibitem{holes}
G.~G\'evay and G.~A.~Jones,
\emph{Hole operations on Hurwitz maps}, Sections 2 and 6,
\url{https://arxiv.org/abs/2111.05566}.
\end{thebibliography}
\end{document}
